\section{Evaluation}
\label{sec:evaluation}

\subsection{Initial Simulation Results}

Our initial experiments demonstrate that the GRU can learn
simulated tree-PLRU replacement behavior with near-perfect
accuracy across all three associativities.

\begin{table}[htbp]
\centering
\caption{GRU Performance on Simulated Tree-PLRU}
\label{tab:gru_results}
\begin{tabular}{|c|c|c|c|}
\hline
\textbf{Ways} & \textbf{Train Samples} &
\textbf{Test Accuracy} & \textbf{Train Time} \\
\hline
4  & 200,000   & 100\%    & 3.3 min  \\
8  & 200,000   & 99.988\% & 3.8 min  \\
16 & 1,000,000 & 100\%    & 12.1 min \\
\hline
\end{tabular}
\end{table}

All models reached the target validation accuracy of
99.9\% within four epochs. The 16-way model required
more training data, reflecting the increased complexity
of learning replacement behavior for larger caches.

These results demonstrate that our model can accurately
predict eviction victims under controlled simulation
conditions. However, the test sequences follow the same
distribution as the training data and contain no
measurement noise, so performance on real hardware
remains to be evaluated.

\subsection{Adaptive Policy Results}

We trained the 16-way DRRIP model for ten epochs, and it
reached 96.1\% test accuracy. As
Table~\ref{tab:drrip_results} shows, accuracy decreased as
the number of misses in a history increased, especially in
SRRIP mode. Each miss moves a new block into the set and
can age every way's counter, so the model must track much
more state than for PLRU. Accuracy also decreased on
histories longer than those used in training, which
contain more misses.

We also trained the same model without the mode token. In
that case identical histories carry different labels, so
no model can exceed about 66\% accuracy. It reached
64.5\%, close to this limit, which confirms that the mode
cannot be inferred from a single set's history and must
be supplied.

To check whether this effect is specific to DRRIP, we
trained the same model on the same histories labeled with
tree-PLRU. It reached 73.6\% test accuracy, and its accuracy
decreased more steeply with misses
(Table~\ref{tab:drrip_results}). With hits only, the PLRU
victim depends only on the order of recent accesses. With
misses, each new block is placed in the way the tree
selects, so the model must track which block occupies each
way as well as the tree bits. Histories without misses were
still predicted with 99.5\% accuracy even at 33--128
accesses, so the model generalizes to longer histories when
blocks do not move.

This could be a problem for our hardware experiments.
Distinguishing adaptive policies requires histories with
misses, and the policies we expect on our hardware targets
also rely on per-way counters. Since accuracy decreased
with misses for both policies, tracking blocks across
misses is the main challenge. In our next steps, we will
investigate whether a larger model, more training histories
with many misses, or predicting the evicted block instead
of the way improve accuracy in these cases.

\begin{table}[htbp]
\centering
\caption{GRU Accuracy on the Same Simulated 16-Way Histories by Misses in the History}
\label{tab:drrip_results}
\begin{tabular}{|c|c|c|c|}
\hline
\textbf{Misses} & \multicolumn{2}{c|}{\textbf{DRRIP}} & \textbf{PLRU} \\
\cline{2-3}
 & \textbf{SRRIP} & \textbf{BRRIP} & \\
\hline
0       & 100\%   & 100\%   & 100\%   \\
3--4    & 98.7\%  & 97.3\%  & 65.1\%  \\
5--8    & 90.0\%  & 94.4\%  & 39.9\%  \\
9--16   & 53.2\%  & 92.6\%  & 20.7\%  \\
\hline
\end{tabular}
\end{table}

\subsection{Planned Model Evaluation}

We plan to evaluate generalization using longer access
histories of up to 512 accesses and diagnostic sequences
where PLRU and LRU predict different eviction victims.
We will also measure per-victim accuracy, model confidence,
training time, and sample efficiency.

\subsection{Hardware Evaluation and Policy Comparison}

After collecting hardware measurements from the L2 of the
Intel Celeron N3450, we will evaluate
the model's ability to predict actual eviction outcomes.
We will compare these predictions against known replacement
policies, including LRU, FIFO, PLRU, and RRIP, by measuring
the percentage of access histories for which each policy
predicts the observed victim.

Diagnostic sequences will help distinguish between
candidate policies. Our model is accurate on hits-only
histories, so we will use it for policy comparison on
those first. Since our DRRIP results show that adaptive
modes are nearly indistinguishable on hits-only histories,
we will add histories with misses where needed, and compare
those directly against the simulators until the model is
accurate on them. If
no known policy consistently
matches hardware behavior, we will investigate whether
the processor uses an adaptive, randomized, or
undocumented replacement mechanism.
